\section{Dataset and Splits}
\label{sec:data}
\subsection{Chinese Dish Data}
Our main source is CNFOOD-241 \cite{cnfood241}, a Chinese dish dataset released
with 241 categories and 191\,811 images at $600\times600$ pixels, together with
the fine-grained recognition study of \cite{xiong2026cnfood}. Part of its imagery
derives from ChineseFoodNet \cite{chen2017chinesefoodnet}; the release is
licensed CC BY-NC 3.0 and is used here for non-commercial research only, and no
image is redistributed with this paper. The release ships \cnPrePoolN{} training and
\cnPreTestN{} validation images, \cnAllN{} in total. We re-partition the training
portion into
\cnPreTrainN{} training and \cnPreValN{} validation images so that the validation
set that drives checkpoint selection is never the one used for reporting, and we
report on the released validation portion as our test split. All counts are exact directory counts and are emitted
into the paper by the build script, so they cannot drift from the code.

\subsection{The Twenty Canteen Dishes and the Catch-All Bucket}
A canteen menu is short. We select twenty dishes that are served in Chinese
campus canteens and that span the cooking methods a tray actually contains:
stir-fried potato shreds, di san xian (potato, aubergine and pepper), hand-torn
cabbage, stir-fried greens, dry-fried green beans, tomato and egg, steamed egg,
sweet-and-sour ribs, cola chicken wings, poached chicken, braised chicken, spicy
chicken, braised pork, chilli pork, sweet-and-sour pork, farmhouse pork, pepper
pork, mu xi pork, baozi and fried noodles. The selection deliberately contains
pairs that a diner distinguishes by one cue only---ribs against pork, braised
against twice-cooked pork, three cuts of pork in the same sauce family---because
those pairs are where a deployed system actually fails.

Twenty classes are not a closed world, and a canteen serves more than twenty
dishes. We therefore add a twenty-first class, \emph{other Chinese dish}, whose
training images are drawn from the 175 categories of CNFOOD-241 that remain after
removing the twenty dishes, seventeen Western dessert classes and twenty-nine
held-out classes reserved for a separate out-of-distribution probe. At most
twelve images are taken from each remaining category, giving 2\,100 training,
350 validation and 1\,396 test images. The bucket is not a junk drawer: it makes
"some other dish" an explicit, trained outcome, so that a confident wrong answer
inside the menu and an honest abstention are different events that can be counted
separately. Table \ref{tab:cndata} lists the resulting splits.

\begin{table}[t]
\centering
\caption{Chinese canteen data. The 21-way task is twenty dishes plus a catch-all
bucket; the 241-class stage uses the full label space. Counts are exact.}
\label{tab:cndata}
\begin{tabular}{llrr}
\toprule
Stage & Split & Classes & Images \\
\midrule
Pre-training & train & 241 & \cnPreTrainN \\
             & val   & 241 & \cnPreValN \\
             & test  & 241 & \cnPreTestN \\
\midrule
Canteen head & train & 21 & \cnTrainN \\
             & val   & 21 & \cnValN \\
             & test  & 21 & \cnTestN \\
\bottomrule
\end{tabular}
\end{table}

\subsection{External Test Sets}
Every number computed on the splits above shares a source with the training data,
so we add two external sets that no part of the pipeline has seen.

\textbf{Web images of the twenty dishes.} We query an image search engine with
each dish name in Chinese, download the first three result pages, and keep
\webN{} images after screening. Screening removes collages, watermarked stock
photography, illustrations, price tags, images whose short side is below 400
pixels, and images whose subject occupies less than a third of the frame. Every
candidate is then compared against a fingerprint of all 24\,889 images in the
21-way training, validation and test splits: an MD5 match or a 64-bit dHash
within Hamming distance four is rejected, which removed four exact and 27 near
duplicates. The remaining images were reviewed one by one; 100 were removed
because the label was wrong or the picture was not a plated dish, and the labels
were corrected before any model was evaluated on them. The set is frozen and its
manifest is part of the released artefacts.

\textbf{Sixteen famous Chinese dishes.} A public Kaggle set of sixteen dishes
\cite{patzer16chinese} overlaps our menu in exactly one class,
sweet-and-sour pork. After the same duplicate screening, \extN{} images remain,
which is enough to sanity-check one class and not enough to characterise the
system; we report it as such.

\subsection{Real Photographs and the Western Baseline}
The project began on a Western menu, and we keep that study as a baseline rather
than deleting it. It uses a twenty-class subset of Food-101
\cite{bossard2014food101} with the official split, and a hundred real
photographs taken in a campus cafeteria of sixteen of those classes
(\texttt{canteen\_real}), which serve as a genuine domain shift for that branch.
The canteen branch of this paper uses no synthetic proxy: its domain shift is the
real one between curated dish photography and hand-held photographs, measured on
the web set above.

\subsection{Energy Priors}
Each of the twenty dishes carries an energy density in kcal per 100\,g and a
typical canteen portion mass, taken from public food-composition tables and
rounded to one representative value per class. The point estimate is their
product; the interval is described in Section \ref{sec:calorie}. These are
population averages, not measurements of the photographed plate, and the
interface states this explicitly.